\section{Results and Discussion}
\label{sec:results}

We evaluate the proposed formulation in cluttered multi-obstacle clearing
tasks with three objectives: to compare its performance against representative
baselines, to assess the contribution of frontier-gap-guided prioritization
and structured contact generation through ablations, and to verify that the
validated interventions remain executable on hardware.

\subsection{Quantitative Evidence}
\label{subsec:quantitative}

All numerical experiments are implemented in Python3 and PyBullet\cite{coumans2019}. The simulation uses a step of $\Delta t=1/240\,\mathrm{s}$ and a control period of $1/40\,\mathrm{s}$. The nominal scenario is an $8\times5\,\mathrm{m}$ bounded workspace in which two robots must clear a passage for the external vehicle through scenes containing $10$--$15$ movable obstacles placed randomly under a minimum separation. Movable obstacles have masses uniformly sampled from $[5,15]\,\mathrm{kg}$, while immovable structures are modeled with zero mass. A trial is counted as successful when a $W$-clear path exists from start to goal and the target reaches its goal disk. For fairness, all baselines are evaluated under the same $W$-clearance criterion and contact model; for the proposed method, each task rollout uses a horizon of $80$ simulation steps, and up to $64$ candidate push tasks are generated per expansion.

% All numerical experiments are implemented in Python3 and PyBullet\cite{coumans2019}. The
% nominal setting uses an $8 \times 5$\,m bounded workspace in which two robots
% must clear a passage for an external vehicle through scenes containing 10--15
% randomly placed movable obstacles. A trial is counted as successful when a
% $W$-clear path exists from start to goal and the target reaches its goal
% disk. For fairness, all baselines are evaluated under the same $W$-clearance
% criterion and contact model.

Table~\ref{tab:main_ablation} summarizes the main quantitative results. We
compare the proposed method against three representative baselines:
\emph{DFS-WCCG}, a simulation-in-the-loop depth-first search guided by the
WCCG with four fixed axis-aligned pushes for each movable obstacle;
\emph{SL-Push}, which follows a straight or waypointed corridor from start to
goal, clears the obstacles intersecting that route, and then validates those
route-induced pushes with physics; and \emph{Rec-NAMO}, a representative
traditional NAMO baseline that incrementally improves local traversability
but often fails to establish a complete $W$-clear corridor for the large
vehicle. Across the evaluated scenes, the proposed method achieves the
highest success rate at 92.5\% while also maintaining the lowest combined
planning and execution time among the successful planners. In contrast,
\emph{DFS-WCCG} succeeds in only 25.0\% of cases because graph-guided
branching quickly becomes intractable. \emph{SL-Push} improves physical
feasibility relative to route-only clearing, but it requires many simulation
calls and longer action sequences. \emph{Rec-NAMO} is faster than
simulation-heavy alternatives, yet it often produces only partial progress
rather than a fully connected $W$-clear passage from start to goal. Taken
together, these comparisons support the need to couple frontier-gap-guided
bottleneck reasoning with contact-feasible action generation.

The ablation study further clarifies which parts of the proposed formulation
support this performance. Removing \texttt{presearch} reduces success from
92.5\% to 85.0\% and increases the number of simulations, suggesting that
frontier-gap-guided prioritization helps reduce unnecessary validation and
focus computation on useful bottlenecks. Removing \texttt{ModeTable} reduces
success further to 80.0\%, again with higher planning cost and more
simulation calls, indicating that structured reuse of contact modes improves
planning efficiency while helping preserve executable action generation under
the same task setting. Together, these ablations support both sides of the
proposed interface: frontier-gap-guided bottleneck prioritization on the
representation side, and structured reuse of contact-feasible action
templates on the execution side.

\subsection{Hardware Execution}
\label{subsec:hardware}


% To verify that the frontier-gap-guided interventions remain valid beyond simulation, we deploy the method on a real two-robot platform in cluttered corridor-construction tasks. Each hardware trial is conducted in a $5\times6\,\mathrm{m}$ workspace and uses $2$ UGVs together with $6$ movable obstacles. Both robots and movables are tracked by motion-capture markers, and the observed poses are streamed to PyBullet for online execution and visualization. Movables are randomly placed at the beginning of each trial to generate diverse clutter. A representative execution is shown in Fig.~\ref{fig:hardware}. Starting from an initially blocked scene, the robots sequentially open a series of frontier gaps until a $W$-clear corridor emerges from start to goal.

To verify that the frontier-gap-guided interventions remain valid beyond
simulation, we deploy the method on a real two-robot platform in cluttered
corridor-construction tasks. Each trial is conducted in a $5 \times 6$\,m
workspace with 2 UGVs and 6 movable obstacles. The scene state is observed
online through motion capture and used to update execution as well as the
current WCCG throughout the task.

A representative execution is shown in Fig.~\ref{fig:hardware}. Starting from
an initially blocked scene, the robots sequentially open frontier gaps until a
$W$-clear corridor emerges from start to goal. As obstacles are displaced, the
WCCG is rebuilt from the observed scene configuration, so the planner
continues reasoning over the currently reachable bottlenecks rather than
following a fixed open-loop obstacle-removal sequence. This is important in hardware because small motion deviations can invalidate a purely geometric plan, whereas the updated graph keeps the intervention logic aligned with the evolving scene.

Overall, the hardware experiments support both sides of the proposed
interface. On the representation side, the updated WCCG continues to expose
the relevant frontier gaps as the scene evolves. On the execution side, the validated gap-opening actions remain executable in the representative hardware runs. Together, these results show that frontier gaps remain a useful
interface between clearance reasoning and contact-feasible multi-robot
pushing in representative real-world executions.


